% Exemplo: pôster (painel) de evento acadêmico, 90 x 120 cm, com tikzposter.
% É o mesmo pôster do exemplo com beamerposter, refeito com a classe
% tikzposter: faixa com o nome do evento, título, autor e seções em duas
% colunas. No tikzposter a página inteira é uma figura TikZ: o fundo é um
% "background style", cada seção é um \block (caixa medida e desenhada pela
% classe) e as colunas são lado a lado, sem quebra automática de uma coluna
% para a outra. Toda a arte é desenhada com TikZ e as figuras são geradas no
% próprio documento: não há imagens de terceiros.
% Compilar: latexmk -pdf poster-tikzposter.tex (biber para as referências)
% Testado com: tikzposter 2.0, biblatex-abnt 3.4,
% TeX Live 2024 ou mais recente
%%BEGIN:tikzposterpreambulo
% 25pt: tamanhos grandes da classe (trocados abaixo); margin=0mm: a arte
% vai até a borda; innermargin: margem do conteúdo; colspace: espaço entre
% colunas
\documentclass[25pt, portrait, margin=0mm, innermargin=6cm, colspace=4.68cm,
  blockverticalspace=1.05cm]{tikzposter}
% tamanho do painel em cm (o evento define): a classe só conhece A0 a A2
\geometry{paperwidth=90cm, paperheight=120cm}
% a classe mede a largura útil antes do \geometry: refaz a conta
\makeatletter
\setlength{\TP@visibletextwidth}{\textwidth-2\TP@innermargin}
\setlength{\TP@visibletextheight}{\textheight-2\TP@innermargin}
\makeatother
\usepackage[T1]{fontenc}
\usepackage[brazilian]{babel}
\usepackage{helvet}         % sem serifa, com as medidas da Arial
\renewcommand{\familydefault}{\sfdefault}
\usepackage[helvet]{sfmath} % matemática sem serifa, como a da beamer
\usepackage[varqu,varl,scaled=0.95]{inconsolata}  % fonte de máquina do livro
% o inconsolata carrega o upquote no \begin{document}, mas o tikzposter
% abre ali a figura da página, antes (erro "Loading a class or package
% in a group"): a regra põe o código do tikzposter depois do código do
% inconsolata
\DeclareHookRule{begindocument}{tikzposter}{after}{inconsolata}
\usetikzlibrary{positioning}
\usepackage{pgfplots}       % gráfico de barras
\pgfplotsset{compat=1.18}
\usepackage{booktabs}
\usepackage{xurl}           % endereços longos quebram em qualquer ponto
\usepackage{ragged2e}       % texto justificado (\justifying)
\usepackage{enumitem}       % recuo e marcador das listas
\usepackage{caption}        % \captionof: legendas fora de flutuantes
\usepackage{csquotes}
% [S. l.] e [s. n.] saem quando faltam local ou editora (NBR 6023:2025)
\usepackage[backend=biber, style=abnt]{biblatex}
\addbibresource{poster-evento.bib}
% Referências: tamanho da letra, entrelinha e espaço entre uma referência
% e a seguinte. Para mudar a aparência da lista, altere só estes três
% valores.
% tamanho da letra (o texto do pôster tem 36 pt)
\newcommand{\refletra}{24pt}
% distância entre as linhas de uma mesma referência
\newcommand{\refentrelinha}{29pt}
% espaço entre as referências; a NBR 6023 pede uma linha em branco
% (\baselineskip), que num pôster ocupa muito espaço
\newcommand{\refespaco}{0pt}
% O biblatex compõe a lista com \bibfont (e não com o tamanho em vigor); o
% biblatex-abnt usa \normalfont no subtítulo, que aqui voltaria ao tamanho
% do texto: na lista, ele só tira o negrito e o itálico. O espaço entre as
% referências entra no \bibsetup: o estilo abnt o redefine ali, ao imprimir
% a lista, e por isso um \setlength comum no preâmbulo seria desfeito.
\renewcommand*{\bibfont}{\fontsize{\refletra}{\refentrelinha}\selectfont
  \renewcommand*{\normalfont}{\mdseries\upshape}}
\AtBeginDocument{\appto{\bibsetup}{\setlength{\bibitemsep}{\refespaco}}}
%%END:tikzposterpreambulo

%%BEGIN:tikzposterestilo
% Tamanhos de letra pedidos pelo evento: texto com 36 pt (legível a 1 m),
% títulos das seções com 40 pt. A classe compõe o texto dos blocos com
% \large e o título dos blocos com \LARGE: os dois são redefinidos aqui.
\renewcommand{\normalsize}{\fontsize{36}{44}\selectfont}
\renewcommand{\small}{\fontsize{30}{36}\selectfont}
\renewcommand{\large}{\fontsize{36}{44}\selectfont}
\renewcommand{\LARGE}{\fontsize{40}{48}\selectfont}
\normalsize
% parágrafos sem recuo, como na beamer (\justifying usa \JustifyingParindent)
\setlength{\parindent}{0pt}
\setlength{\JustifyingParindent}{0pt}

% Cores do painel
\definecolor{eventoescuro}{HTML}{0B3C5D}
\definecolor{eventomedio}{HTML}{1D7A8C}
\definecolor{eventodestaque}{HTML}{E07A1F}
\definecolor{eventotexto}{HTML}{4A4A4A}
\colorlet{titlefgcolor}{eventoescuro}
\colorlet{blocktitlefgcolor}{eventoescuro}
\colorlet{blockbodyfgcolor}{eventotexto}
\tikzposterlatexaffectionproofoff   % sem a marca "LaTeX TikZposter" no canto

% Arte do painel, desenhada com TikZ: faixa do alto com o nome do evento,
% ícone de livro aberto e círculos; ondas no pé. No tikzposter a origem
% fica no centro da página: o scope leva a origem ao canto de baixo, à
% esquerda.
\newcommand{\nomeevento}{SIMPÓSIO DE EXEMPLO 2026}
\newcommand{\lemaevento}{Ensino, pesquisa e extensão}
\definebackgroundstyle{evento}{%
  \begin{scope}[shift={(bottomleft)}, x=1cm, y=1cm]
    \fill[white] (0,0) rectangle (90,120);
    % faixa do alto, com a borda de baixo em onda
    \shade[left color=eventoescuro, right color=eventomedio]
      (0,120) -- (90,120) -- (90,106.5)
      .. controls (65,102.5) and (35,109) .. (0,104.5) -- cycle;
    % círculos translúcidos
    \foreach \x/\y/\r in {80/114/9, 69/119/5, 86/104/4, 8/118/4, 58/108/2.5}
      \fill[white, opacity=.07] (\x,\y) circle (\r);
    % ícone: livro aberto
    \begin{scope}[shift={(13.5,112.3)}, line width=4pt, white]
      \draw[fill=white, fill opacity=.15]
        (0,0) .. controls (-2.5,1) and (-5,1) .. (-7,0.2)
        -- (-7,-6.8) .. controls (-5,-6) and (-2.5,-6) .. (0,-7) -- cycle;
      \draw[fill=white, fill opacity=.15]
        (0,0) .. controls (2.5,1) and (5,1) .. (7,0.2)
        -- (7,-6.8) .. controls (5,-6) and (2.5,-6) .. (0,-7) -- cycle;
      \foreach \k in {1,2,3} {
        \draw[line width=2.5pt] (-5.8,-1.2*\k)
          .. controls (-4,-1.2*\k+.4) .. (-1.2,-1.2*\k-.2);
        \draw[line width=2.5pt] (5.8,-1.2*\k)
          .. controls (4,-1.2*\k+.4) .. (1.2,-1.2*\k-.2);
      }
    \end{scope}
    \node[anchor=west, text=white, font=\fontsize{80}{90}\selectfont\bfseries]
      at (25,113.2) {\nomeevento};
    \node[anchor=west, text=white, font=\fontsize{40}{48}\selectfont]
      at (25.3,108.6) {\lemaevento};
    % fio de destaque abaixo da faixa
    \draw[eventodestaque, line width=6pt] (6,101.5) -- (84,101.5);
    % ondas do pé
    \fill[eventomedio, opacity=.85] (0,0) -- (0,4.5)
      .. controls (25,8.5) and (55,1) .. (90,5.5) -- (90,0) -- cycle;
    \fill[eventodestaque, opacity=.9] (0,0) -- (0,2.2)
      .. controls (30,0.2) and (60,5.5) .. (90,2.5) -- (90,0) -- cycle;
  \end{scope}}
\usebackgroundstyle{evento}

% Título do trabalho, autor e instituição, sem caixa em volta (o estilo do
% título não desenha nada; as chaves dão a largura e as distâncias). A
% classe guarda o título numa \scalebox, onde \\ não quebra a linha: \title
% é redefinido para guardar o texto como está.
\makeatletter
\renewcommand{\title}[1]{\gdef\@title{#1}}
% a chave titletextscale da classe escreve "1=1" na página (sai invisível,
% mas o log acusa "Missing character"): a chave passa a guardar só o valor
\define@key{title}{titletextscale}{\def\TP@titletextscale{#1}}
\definetitlestyle{evento}{width=78cm, innersep=0pt, linewidth=0pt,
  titletotopverticalspace=19.05cm, titletoblockverticalspace=2.5cm}{}
\usetitlestyle{evento}
% composição do título (o ambiente settitle da classe faz o mesmo, mas a
% definição fica presa no grupo do ambiente)
\renewcommand{\TP@maketitle}{%
  \centering
  {\fontsize{72}{84}\selectfont\bfseries\color{eventoescuro}\@title\par}
  \vspace{1.5cm}
  {\LARGE\color{eventotexto}\@author\par\@institute\par}}
\makeatother

% Título de seção: caixa arredondada com contorno na cor do evento. A
% classe mede o título e o texto e passa os nós blocktitle e blockbody ao
% estilo, que só desenha a moldura do título (o texto não tem caixa).
\defineblockstyle{evento}{titlewidthscale=1, bodywidthscale=1, titlecenter,
  titleoffsetx=0pt, titleoffsety=0pt, bodyoffsetx=0pt, bodyoffsety=0pt,
  bodyverticalshift=.7cm, roundedcorners=0, linewidth=0pt,
  titleinnersep=0pt, bodyinnersep=0pt}{%
  \ifBlockHasTitle
    \draw[eventomedio, line width=2.25pt, rounded corners=.65cm]
      ([shift={(1.125pt,-1.125pt)}]blocktitle.north west) rectangle
      ([shift={(-1.125pt,1.125pt)}]blocktitle.south east);
  \fi}
\useblockstyle{evento}
% \secao[chaves]{título}{texto}: bloco com o título em maiúsculas, numa
% caixa de 3,3 cm de altura (a régua invisível fixa a altura e centra as
% maiúsculas), e o texto justificado, com hifenização. A caixa vazia
% (\null) no começo do texto faz a primeira linha ficar a uma entrelinha
% dela, como as seguintes (sem ela, a distância até o título mudaria com o
% tamanho da letra)
\newcommand{\secao}[3][]{%
  \block[#1]{\rule[-1.14cm]{0pt}{3.3cm}\MakeUppercase{#2}}%
    {\null\justifying #3}}

% Legendas numeradas acima da figura e da tabela (como pede a ABNT); o
% tikzposter não tem flutuantes (o bloco é uma caixa dentro da figura TikZ
% da página): o ambiente ilustracao centraliza e \captionof numera a legenda.
% \fonte escreve a fonte abaixo
\captionsetup{font=small, labelfont=bf, labelsep=colon, skip=0pt}
% [#1] espaço depois (0pt na última linha de um bloco), #2 figure ou table,
% #3 título, #4 rótulo
\newenvironment{ilustracao}[4][.4cm]{%
  \def\depoisdailustracao{#1}%
  \par\vspace{.9cm}\begingroup\centering
  \captionof{#2}{#3}\label{#4}\par\vspace{.1cm}}%
  {\par\endgroup\vspace{\depoisdailustracao}}
\newcommand{\fonte}[1]{\par\vspace{.4cm}{\small Fonte: #1\par}\vspace{1cm}}
% listas (no bloco, o topsep é ignorado)
\setlist[itemize]{label=\textcolor{eventodestaque}{\textbullet},
  leftmargin=1.5cm, labelsep=.5cm, topsep=0pt, itemsep=.1cm, parsep=0pt,
  before=\vspace{.1cm}, after=\vspace{.1cm}}
% a lista de referências começa logo abaixo do título da seção
\AtBeginBibliography{\setlength{\topsep}{0pt}\setlength{\partopsep}{0pt}}
% \cmd{nome}: comando do LaTeX destacado no código da Figura 2
\newcommand{\cmd}[1]{\textcolor{eventomedio}{\bfseries\textbackslash #1}}
% \marca{n}: marcador numerado, como os do livro (círculo laranja)
\newcommand{\marca}[1]{\tikz[baseline=(m.base)]%
  \node[circle, fill=eventodestaque,
  text=white, inner sep=.12cm, font=\small\bfseries] (m) {#1};}
%%END:tikzposterestilo

%%BEGIN:tikzposterdados
% título em duas linhas (72 pt): \\ marca onde a linha quebra
\title{LaTeX para documentos acadêmicos:\\ um livro e um guia aberto de
  exemplos testados}
\author{Rodrigo Smarzaro}
\institute{Departamento de Exemplo, Universidade Exemplo
  (autor@exemplo.edu.br)}
% Apoio financeiro, centralizado no pé em coluna única: troque as imagens
% de exemplo (pacote mwe) pelos logotipos dos financiadores; \apoiofalse
% retira a seção
\newif\ifapoio
\apoiotrue
\newcommand{\apoiofinanceiro}{%
  % .45 da largura útil; offsets: sobe o bloco, mais perto das colunas
  \secao[titlewidthscale=.45, bodywidthscale=.45,
    titleoffsety=.9cm, bodyoffsety=.9cm]{Apoio financeiro}{%
    \centering
    \includegraphics[height=3cm]{example-image-a}\hspace{2cm}%
    \includegraphics[height=3cm]{example-image-b}\hspace{2cm}%
    \includegraphics[height=3cm]{example-image-c}}}
%%END:tikzposterdados

\begin{document}
%%BEGIN:tikzpostertitulo
\maketitle
%%END:tikzpostertitulo

%%BEGIN:tikzpostercolunas
% \column{.5}: duas colunas iguais; cada \block fica abaixo do anterior
\begin{columns}
\column{.5}
  \secao{Introdução}{%
    Trabalhos acadêmicos seguem normas detalhadas de formatação. O \LaTeX,
    criado por Lamport sobre o \TeX\ de Knuth, separa o conteúdo da forma:
    o autor escreve o texto e marca a estrutura, e o programa cuida da
    numeração, das referências cruzadas, das citações e da lista de
    referências \parencite{Lamport1994}. Mesmo assim, quem começa encontra
    material em português disperso, muitas vezes desatualizado, e exemplos
    que não compilam mais.}

  \secao{Objetivos}{%
    \begin{itemize}
      \item Oferecer um livro em português sobre \LaTeX\ para trabalhos
        acadêmicos, artigos, apresentações e material didático.
      \item Manter um guia público, com licença aberta, em que todos os
        exemplos do livro compilam \parencite{latexacademico}.
      \item Mostrar o uso das normas da ABNT do início ao fim.
    \end{itemize}}

  \secao{Metodologia}{%
    Cada exemplo nasce no guia e é compilado por testes automáticos a cada
    versão, com o \TeX\ Live mais recente. O livro não copia o código:
    inclui o trecho do próprio arquivo do guia, delimitado por marcadores, e
    a figura do resultado gerada a partir dele (Figura~\ref{fig:fluxo}).
    \begin{ilustracao}{figure}{Do guia público ao livro}{fig:fluxo}
      \begin{tikzpicture}[caixa/.style={draw=eventomedio, line width=2pt,
          rounded corners=.5cm, minimum width=13cm, minimum height=4.5cm,
          align=center}]
        \node[caixa] (guia) {\textbf{Guia público}\\exemplos testados};
        \node[caixa, right=7cm of guia] (livro) {\textbf{Livro}\\texto e
          figuras};
        \draw[-latex, line width=4pt, eventodestaque] (guia) -- (livro)
          node[midway, above=.3cm, font=\small] {trechos e}
          node[midway, below=.3cm, font=\small] {resultados};
      \end{tikzpicture}
      \fonte{elaborado pelo autor.}
    \end{ilustracao}
    No livro, cada exemplo forma uma unidade com o código e o resultado lado
    a lado, e marcadores numerados ligam o texto às linhas do código: a
    seção \marca{1} é numerada sozinha, e a fórmula \marca{2} sai em
    destaque, centralizada (Figura~\ref{fig:unidade}).
    \begin{ilustracao}[0pt]{figure}
      {Uma unidade Exemplo: código e resultado}{fig:unidade}
      \begin{minipage}[c]{.56\linewidth}
        \begin{tikzpicture}
          \node[fill=black!5, draw=black!25, line width=1.5pt, inner sep=.8cm,
                font=\ttfamily\small, align=left] {%
            \begin{tabular}{@{}r@{\hspace{.5cm}}l@{}}
              \smash{\marca{1}} & \cmd{section}\{Introdução\} \\
                                & A média de \$n\$ valores é \\
              \smash{\marca{2}} & \cmd{[} \cmd{bar}\{x\} =
                \cmd{frac}\{1\}\{n\} \\
                                & \hspace*{1.5em}\cmd{sum}\_\{i=1\}%
                                  \textasciicircum\{n\} x\_i \cmd{]}
            \end{tabular}};
        \end{tikzpicture}
      \end{minipage}\hfill
      \begin{minipage}[c]{.4\linewidth}
        \begin{tikzpicture}
          \node[fill=white, draw=black!25, line width=1.5pt, inner sep=.8cm,
                align=left, text width=12.3cm, font=\rmfamily\small] {%
            {\bfseries 1\quad Introdução}\par\medskip
            A média de $n$ valores é
            \[ \bar{x} = \frac{1}{n}\sum_{i=1}^{n} x_i \]};
        \end{tikzpicture}
      \end{minipage}
      \fonte{elaborado pelo autor.}
    \end{ilustracao}}

\column{.5}
  \secao{Resultados}{%
    O livro e o guia cobrem os principais tipos de documento acadêmico
    (Tabela~\ref{tab:numeros}), e os exemplos de trabalhos acadêmicos
    seguem a NBR 14724 \parencite{NBR14724}, com a classe abnTeX2
    \parencite{abntex2}. A Figura~\ref{fig:temas} mostra os temas com mais
    arquivos de exemplo no guia.\par\vspace{.3cm}
    % tabela e figura lado a lado, alinhadas pela primeira linha da legenda
    \begin{minipage}[t]{.46\linewidth}
      \begin{ilustracao}[0pt]{table}{Livro e guia em números}{tab:numeros}
        \begin{tabular}{@{}lr@{}}
          \toprule
          \textbf{Item}          & \textbf{Total} \\
          \midrule
          Capítulos              & 17 \\
          Apêndices              & 4 \\
          Exemplos comentados    & 81 \\
          Trechos de código      & 164 \\
          Exemplos no guia       & 88 \\
          Receitas no guia       & 81 \\
          \bottomrule
        \end{tabular}
        \fonte{elaborado pelo autor.}
      \end{ilustracao}
    \end{minipage}\hfill
    \begin{minipage}[t]{.52\linewidth}
      \begin{ilustracao}[0pt]{figure}{Exemplos por tema no guia}{fig:temas}
        \begin{tikzpicture}
          \begin{axis}[xbar, width=12cm, height=15.5cm, xmin=0, xmax=12,
              bar width=.9cm, symbolic y coords={Áreas, Figuras, Tabelas,
                Apresentações, Diversão, Referências, Básico, Código},
              ytick=data, axis x line=none, axis y line*=left,
              nodes near coords={\pgfmathprintnumber[assume math mode=true]%
                {\pgfplotspointmeta}}, % números sem $...$: na fonte do texto
              every node near coord/.style={font=\small},
              tick label style={font=\small}]
            \addplot[fill=eventomedio!75, draw=none] coordinates {(5,Áreas)
              (6,Figuras) (6,Tabelas) (8,Apresentações) (8,Diversão)
              (8,Referências) (10,Básico) (10,Código)};
          \end{axis}
        \end{tikzpicture}
        \fonte{elaborado pelo autor.}
      \end{ilustracao}
    \end{minipage}}

  \secao{Conclusões}{%
    Com os exemplos testados a cada versão, o código do livro continua
    compilando quando os pacotes mudam \parencite{Mittelbach2023}. O guia
    está aberto a pedidos de exemplos e correções, e o leitor encontra nele
    o arquivo completo de cada exemplo, pronto para compilar.}

  % referências no fim da segunda coluna (só o apoio financeiro ocupa a
  % largura toda)
  \secao{Referências}{\printbibliography[heading=none]}
\end{columns}
%%END:tikzpostercolunas

%%BEGIN:tikzposterapoio
\ifapoio\apoiofinanceiro\fi
%%END:tikzposterapoio

\end{document}
